% !PRINTED 279
\weberpage{279}
%<<CONT
for from the expansion in this one case one can derive the
expansions for the remaining cases, by replacing $\omega$
by
\weq{17}{\dfrac{\omega}{n}, \quad \dfrac{\omega}{n'}, \quad \dfrac{\omega}{n''}}
and then adding integers to $\omega$, whereby no negative
powers of $q$ can be newly introduced.

For carrying out the calculation one makes use of the expansions
\weq{18}{
\begin{array}{l}
U = q^{-\frac{(n'+1)(n''+1)}{24}} \times
\Pi(1+q^{2h-1})(1+q^{(2h-1)n'})(1+q^{(2h-1)n''})(1+q^{(2h-1)n}) \\
U_1 = q^{-\frac{(n'+1)(n''+1)}{24}} \times
\Pi(1-q^{2h-1})(1-q^{(2h-1)n'})(1-q^{(2h-1)n''})(1-q^{(2h-1)n}) \\
U_2 = 4q^{\frac{(n'+1)(n''+1)}{12}} \times
\Pi(1+q^{2h})(1+q^{2hn'})(1+q^{2hn''})(1+q^{2hn}),
\end{array}
}
which in the individual cases furnish the power-series expansions of $A$, $B$,
from which the negative powers of $q$ are to be eliminated. In this way
one computes very simply the following equations:
\weq{19}{
\begin{array}{l}
n = 15, \quad A = 1, \\
n = 21, \quad (A^2-B)^2-A = 0, \\
n = 33, \quad A^2-B-A = 4, \\
n = 35, \quad A^2-B-A = 2, \\
n = 55, \quad A^3-B-4A^2-A+4 = 0.
\end{array}
}

2. If
\[
(n'-1)(n''-1) = 8\mu \equiv 0 \ (\mathrm{mod}\ 8),
\]
then we set
\weq{20}{
\begin{array}{l}
A = \dfrac{uv}{v'v''} + (-1)^\mu\left(\dfrac{u_1v_1}{v_1'v_1''} + \dfrac{u_2v_2}{v_2'v_2''}\right), \\[2ex]
B = \dfrac{v'v''}{uv} + (-1)^\mu\left(\dfrac{v_1'v_1''}{u_1v_1} + \dfrac{v_2'v_2''}{u_2v_2}\right),
\end{array}
}
and obtain, according to (12), the permutations belonging together:
\weq{21}{
\begin{array}{lll}
\omega, & A, & B, \\
-\dfrac{1}{\omega}, & \varrho^{1-n'}A, & \varrho^{n'-1}B, \\
\omega+1, & \sigma^{1-n'}e^{\frac{2\mu\pi i}{3}}A, & \sigma^{n'-1}e^{-\frac{2\mu\pi i}{3}}B.
\end{array}
}
